\chapter{Agents and models}

\label{ch:agents}

This chapter answers two questions that look like one: which model answers, and which program
is asked. They are decided in different places and by different mechanisms. The model comes
from a five-level precedence that ends in the daemon's configuration; the \emph{agent} ---
\texttt{pi} or \texttt{claude} --- comes from the graph, and its command line is assembled
explicitly, flag by flag, because every flag in it was added after its absence produced a
silent failure.

The chapter also states a limit plainly. The role table that would let a graph say
\texttt{model = "@workhorse"} is designed and specified (spec 362) and is \emph{not} in the
tree at this revision; what exists is a node's explicit model, a node's \texttt{tier}, the
run's model, \texttt{\$PII\_MODEL}, and \texttt{[core].model}. A reader who needs the role
table should read the last section of this chapter as a design, not as behaviour.

\section{The five levels that decide a model}

\texttt{crates/pii-server/src/tiers.rs} holds the only written-down statement of the order:

\begin{center}
\resizebox{\linewidth}{!}{%
\begin{tikzpicture}[node distance=8mm and 12mm]
	\node[b] (node) {node \texttt{model}};
	\node[b, below=of node] (tier) {node \texttt{tier}};
	\node[b, below=of tier] (run) {run model\\\scriptsize \texttt{--model}};
	\node[b, below=of run] (env) {\texttt{\$PII\_MODEL}};
	\node[b, below=of env] (cfg) {\texttt{[core].model}};
	\node[v, right=of tier, xshift=16mm] (found) {first set?};
	\node[b, right=of found, xshift=6mm] (use) {use it};
	\node[b, below=of use] (none) {no model\\\scriptsize exit 5};

	\draw[f] (node) -- (tier);
	\draw[f] (tier) -- (run);
	\draw[f] (run) -- (env);
	\draw[f] (env) -- (cfg);
	\draw[f] (tier.east) -- (found.west);
	\draw[f] (found) -- node[p, above=2pt, midway]{yes} (use);
	\draw[sf] (found) -- node[p, above=2pt, midway]{no} (none);
\end{tikzpicture}}
\captionof{figure}{Node model beats node tier beats run model beats environment beats config,
and the order lives in one function so a caller cannot invent a sixth level in its own file.}
\end{center}

Two consequences follow. A graph that names a model explicitly keeps working exactly as it did
before tiers existed --- which is why the precedence has node model at the top rather than the
bottom. And a daemon with no model anywhere refuses to start with exit \texttt{5} rather than
starting and failing at the first dispatch: the failure belongs where the operator is
watching.

\paragraph{Tiers.} A \texttt{[[model\_tier]]} entry names a model class a node can ask for by
name. Each entry has a name, a provider spec, an optional tool allowlist, and an optional
\texttt{max\_tokens} cap. The name must be non-empty and unique, and both are checked at parse
time, because a duplicate would make a node's \texttt{tier} resolve to whichever entry won and
the operator would have no way to know which.

The tool allowlist is enforced by the capability layer, not by the prompt. That distinction is
the whole point of the field: a tier that may call three tools cannot be talked into calling a
fourth, and no instruction in a brief can widen it, because the check happens where the tool
is dispatched. The allowlist only ever shrinks: intersecting two sets takes the names present
in both.

\section{Agent argv, assembled by hand}

The worker node's spec is read from graph data: \texttt{agent}, \texttt{model}, \texttt{spec},
\texttt{workdir}, and optionally \texttt{provider}, \texttt{sid}, \texttt{session},
\texttt{run}, \texttt{max\_budget\_usd}. \texttt{WorkerSpec::from\_cfg} validates each
optional field before anything is spawned: \texttt{max\_budget\_usd} must be a positive finite
number (a typo can never become ``no cap'' or a flag the CLI rejects after the dispatch
started), a claude \texttt{sid} must be a lowercase UUID v4, and a \texttt{run} must match
\texttt{[A-Za-z0-9.\_-]+} because it names a file.

\begin{center}
\resizebox{\linewidth}{!}{%
\begin{tikzpicture}[node distance=8mm and 12mm]
	\node[b] (row) {queue row};
	\node[b, right=of row] (node) {coder node\\\scriptsize model $\cdot$ spec $\cdot$ workdir};
	\node[b, right=of node] (plan) {\texttt{plan()}\\\scriptsize credentials checked};
	\node[b, right=of plan, yshift=14mm] (launch) {launcher script\\\scriptsize sets PATH + env};
	\node[b, right=of plan, yshift=-2mm] (argv) {argv};
	\node[b, right=of argv, yshift=6mm] (agent) {agent process};
	\node[b, below=of plan] (refuse) {refusal\\\scriptsize worker.start.refused};

	\draw[f] (row) -- (node);
	\draw[f] (node) -- (plan);
	\draw[f] (plan) -- (argv);
	\draw[f] (launch) -- (agent);
	\draw[f] (argv) -- (agent);
	\draw[bk] (plan) -- (refuse);
\end{tikzpicture}}
\captionof{figure}{Planning is pure and never spawns, and the spawn goes through a named
launcher script that sets the environment inside itself rather than inheriting the caller's.}
\end{center}

For \texttt{pi}, the argv is built in this order and for these reasons:

\begin{center}
\resizebox{\linewidth}{!}{%
\begin{tikzpicture}[node distance=8mm and 12mm]
	\node[b] (a) {\texttt{pi}};
	\node[b, right=of a] (b) {\texttt{--session}\\\scriptsize file};
	\node[b, right=of b] (c) {\texttt{--model}\\\scriptsize id};
	\node[b, right=of c] (d) {\texttt{--mode json}\\\scriptsize route is readable};
	\node[b, right=of d] (e) {\texttt{--approval-mode}\\\scriptsize write};
	\node[b, below=of e] (f) {\texttt{--provider}\\\scriptsize optional};
	\node[b, left=of f] (g) {\texttt{-p}\\\scriptsize spec + brief};
\end{tikzpicture}}
\captionof{figure}{Every element is a decision that has a failure behind it, not a default
that was copied from a shell history.}
\end{center}

\begin{longtable}{@{}p{0.20\textwidth}p{0.30\textwidth}p{0.42\textwidth}@{}}
\toprule
\textbf{Flag} & \textbf{Value} & \textbf{Why it is there} \\
\midrule
\endhead
\texttt{--session} & \texttt{<workdir>/.pii-sessions/<stem>-<run>.jsonl}, or the node's
\texttt{session} & a headless run needs a session file, and its path is what makes a resume
possible \\
\texttt{--model} & the resolved model id & the only place the model is stated \\
\texttt{--mode json} & \texttt{json} & the event stream names the provider and model per message, so
the run's real route can be read back \\
\texttt{--approval-mode} & \texttt{write} & headless runs have no approval surface; see below \\
\texttt{--provider} & the node's provider, when set & pins the route; the key is \emph{not} on
argv \\
\texttt{-p} & \texttt{Spec: <path>} and the brief & the brief says ``the spec named above'', so the
spec must be named \\
\bottomrule
\end{longtable}

\begin{lstlisting}[language=bash]
# what a pi dispatch looks like, assembled by plan()
pi --session /repo/.pii-sessions/376-retention-r000000000412.jsonl \
   --model deepseek-v4.1-flash:cloud \
   --mode json \
   --approval-mode write \
   --provider ollama-cloud \
   -p "Spec: specs/376-retention-drops-what-the-policy-expires.md

Implement the spec named above. Follow it as the contract. Do not commit."
\end{lstlisting}

For \texttt{claude} the argv is shorter and different in kind:

\begin{lstlisting}[language=bash]
claude --session-id 4f1c2a9e-8b7d-4c31-9a02-6d5e0f7b1c88 \
       --model claude-haiku-5-5 \
       --permission-mode bypassPermissions \
       --output-format stream-json --verbose \
       --max-budget-usd 0.25 \
       -p "<brief>"
\end{lstlisting}

Four decisions are visible in that line. \texttt{--output-format stream-json} with
\texttt{--verbose} is what makes the session id, the cost and the model readable after the
fact; partial messages are \emph{not} requested, because one measured run produced 3\,779
\texttt{stream\_event} lines and passed the capture bound. A resumed run carries
\texttt{--resume <sid>}; a fresh run carries \texttt{--session-id <new>} when the caller chose
one, so the handle exists before the spawn. The two flags are never both emitted. And no
provider flag and no key appear at all: claude reads its own credential.

\section{Approval modes, and what they are not}

Approval modes separate \emph{writing} from \emph{judging}. They are not the security
boundary, and the chapter states that first because it is the mistake to make.

\texttt{pi}'s default is \texttt{always-ask}. Headless, there is no surface to ask on, so
every mutating tool is denied. The measured result was a coder that ran for 122 seconds,
exited 3, and reported:

\begin{lstlisting}[language=bash]
approval required for `bash` but this session has no approval surface
\end{lstlisting}

A coder in that state has written nothing and has reported a plausible reason. The fix was
\texttt{--approval-mode write}, which approves file edits and still refuses to run commands.
That is the right split for this architecture: the coder writes, the judge runs the exit
criteria, and the daemon's capability set stays the outer bound. A coder that could run
arbitrary commands would be running the very checks that are supposed to judge it.

\texttt{claude} has the same problem with a different flag. Headless \texttt{-p} cannot ask
for permission, so without a mode it cannot write at all; the dispatch passes
\texttt{--permission-mode bypassPermissions}. The name is alarming and the reason it is
acceptable is the same: the process is inside the sandbox, under a capability set that decided
what it may do before it started, and its work is judged by criteria it cannot edit.

\begin{longtable}{@{}p{0.14\textwidth}p{0.20\textwidth}p{0.16\textwidth}p{0.40\textwidth}@{}}
\toprule
\textbf{Agent} & \textbf{Mode used} & \textbf{Can write} & \textbf{What still stops it} \\
\midrule
\endhead
\texttt{pi} & \texttt{write} & files only & the capability set, the sandbox, and the judge \\
\texttt{pi} & \texttt{always-ask} (default) & nothing & nothing needs to: the run does nothing and
exits 3 \\
\texttt{claude} & \texttt{bypassPermissions} & yes & the capability set, the sandbox, the judge, and
\texttt{--max-budget-usd} \\
\bottomrule
\end{longtable}

\section{Credentials are checked before the spawn}

A dispatch that names a provider must have that provider's key in its environment. The table
is two entries --- \texttt{ollama-cloud} reads \texttt{OLLAMA\_API\_KEY}, \texttt{deepseek}
reads \texttt{DEEPSEEK\_API\_KEY} --- and a provider not in the table has no key variable, so
the check fails and the dispatch is refused. That is deliberate: an unknown provider is
refused as a missing credential rather than attempted, because attempting it costs a spawn and
produces a provider error instead of a configuration error.

The refusal is journalled as \texttt{worker.start.refused} with the agent, the model and the
reason, and the node fails without a process. This is the difference between a run that fails
in one second with a nameable cause and a run that fails in ninety seconds with a stack trace
from a subprocess.

The key itself never reaches argv. For \texttt{pi}, the value is resolved at runtime from an
environment reference in a pii-owned \texttt{models.json} (\texttt{apiKey = "\$ENV:<NAME>"}),
with \texttt{PI\_CODING\_AGENT\_DIR} pointing at that directory. The measured proof is a run
that exited 0 through ollama-cloud with no \texttt{--api-key} on the command line. An argv key
would be visible in \texttt{/proc/<pid>/cmdline} to every process on the host.

For \texttt{claude}, the credential check accepts either an environment token
(\texttt{ANTHROPIC\_API\_KEY} or \texttt{CLAUDE\_CODE\_OAUTH\_TOKEN}) or the host's login file
with an OAuth block. Only the block's \emph{presence} is checked; its values are never read
into the journal.

\section{The brief must name the spec}

The brief is the text handed to the agent, and the code prefixes it with the spec path:

\begin{lstlisting}[basicstyle=\ttfamily\footnotesize]
let prompt = if spec.spec.is_empty() {
    brief.to_owned()
} else {
    format!("Spec: {}\n\n{brief}", spec.spec)
};
\end{lstlisting}

That two-line prefix exists because of a measured failure. A coder was spawned with the spec
in its session \emph{name} and nothing about the spec in its prompt, and it implemented a
different spec's files: the row said one thing, the tree changed another, and the criteria
that failed belonged to work nobody had asked for. The failure is invisible in the argv log
unless the prompt is logged, which is why the prefix is in the code and pinned by a test
rather than left to a brief author's memory.

\paragraph{An empty brief is refused.} A whitespace-only brief runs the agent with no prompt;
it exits 0 at once, and a criterion that already passes makes the node succeed having done
nothing. The coder refuses such a brief before it spawns. This is the general shape of a guard
that exists because a passing check can be worse than a failing one: success-with-no-work is
the hardest failure to notice.

\section{Two verdicts, and only one of them counts}

The agent's report is a self-report. A coder can exit 0 having written nothing, or having
written the wrong thing, or having weakened a test until it passed. So the node has two inputs
and one rule.

\begin{center}
\resizebox{\linewidth}{!}{%
\begin{tikzpicture}[node distance=8mm and 12mm]
	\node[b] (worker) {agent\\\scriptsize exit code + report};
	\node[b, right=of worker] (verdict) {\texttt{decide()}\\\scriptsize rc, criteria};
	\node[b, above=of verdict] (judge) {judge\\\scriptsize spec blob at a commit};
	\node[v, below=of verdict] (ok) {pass?};
	\node[b, right=of ok, yshift=6mm] (pass) {ok port\\\scriptsize verdict = Passed};
	\node[b, right=of ok, yshift=-6mm] (fail) {err port\\\scriptsize names the criterion};

	\draw[f] (worker) -- (verdict);
	\draw[f] (judge) -- (verdict);
	\draw[f] (verdict) -- (ok);
	\draw[f] (ok) -- node[p, above=2pt, midway]{yes} (pass);
	\draw[bk] (ok) -- node[p, below=2pt, midway]{no} (fail);
\end{tikzpicture}}
\captionof{figure}{The worker decides nothing: its exit code is one input, the criteria runs
are the other, and any non-zero criterion fails the node even when the worker exited cleanly.}
\end{center}

\texttt{decide} has three rules and no exceptions:

\begin{enumerate}
	\item A non-zero worker exit is a failure, named \texttt{worker exited <n>}.
	\item Zero criteria run is a failure, named \texttt{no criteria were run}. This is the rule
	that catches a node whose criteria list was silently empty.
	\item The first criterion with a non-zero exit is a failure, and the message carries the
	criterion command. Every criterion still runs, so the verdict can name more than one
	problem in the log even though the message names the first.
\end{enumerate}

\paragraph{Where the criteria come from.} The spec-mode judge reads them from the spec's blob
\emph{at a commit}, not from the working tree, the run input, or the agent. A worker that
edits its own spec gains nothing: the judge asks git for the document that was reviewed,
parses the \texttt{\#\# Exit Criteria} section, takes the first backticked span of every
checkbox line, and runs each command in the node's workdir. A base that is not exactly forty
lowercase hex digits is refused before git is asked, so a malformed handle is a refusal and
not a mystery.

\begin{center}
\resizebox{\linewidth}{!}{%
\begin{tikzpicture}[node distance=8mm and 12mm]
	\node[b] (commit) {commit\\\scriptsize 40 hex};
	\node[b, right=of commit] (blob) {spec blob};
	\node[b, right=of blob] (parse) {parse criteria};
	\node[b, below=of parse] (each) {run each\\\scriptsize in workdir};
	\node[b, left=of each] (ports) {ok / fail};
	\node[b, left=of ports] (row) {row status\\\scriptsize via record node};

	\draw[f] (commit) -- (blob);
	\draw[f] (blob) -- (parse);
	\draw[f] (parse) -- (each);
	\draw[f] (each) -- (ports);
	\draw[f] (ports) -- (row);
	\draw[sf] (row.north) -- ++(0,10mm) node[p, above=2pt, midway]{a failure does not stop the rest};
\end{tikzpicture}}
\captionof{figure}{The document is read at a commit, so the verdict is about the spec that was
reviewed rather than the tree the agent left.}
\end{center}

\section{The ladder, and when to climb it}

A heal ladder is an ordered list of rungs, cheapest first, each a pair of agent and model.
\texttt{escalate} returns rung zero with no history, and otherwise the next rung only when
\emph{every} previous attempt failed \emph{with evidence}. An attempt that passed, or that
failed without naming a failing criterion, ends the ladder.

The rule sounds fussy and is the difference between escalation and retrying. Evidence is a
failing criterion: a specific command that returned non-zero. Without it there is nothing to
escalate \emph{about} --- the next model would be handed the same problem statement and no new
information, and the spend would buy a second opinion on an unknown question. With it, the
next rung starts from the command that failed.

\begin{center}
\resizebox{\linewidth}{!}{%
\begin{tikzpicture}[node distance=8mm and 12mm]
	\node[b] (try) {attempt n};
	\node[v, right=of try] (ev) {evidence?};
	\node[b, right=of ev, yshift=7mm] (next) {next rung};
	\node[b, right=of ev, yshift=-7mm] (stop) {end ladder};
	\node[b, below=of stop] (cost) {no further spend};

	\draw[f] (try) -- (ev);
	\draw[f] (ev) -- node[p, above=2pt, midway]{criterion failed} (next);
	\draw[f] (ev) -- node[p, above=2pt, midway]{passed or none} (stop);
	\draw[sf] (stop) -- (cost);
	\draw[f] (next.south) |- (try.east);
\end{tikzpicture}}
\captionof{figure}{Escalation is gated on evidence, so a ladder cannot become a slot machine
and a ladder whose rungs all lack evidence never climbs at all.}
\end{center}

\section{The session is the resume handle}

A dispatch records a session before it can need one. For \texttt{pi}, that means writing a
session header first, because an empty or missing session file is rejected before any model is
called. The header is the shape the installed \texttt{pi} emits --- \texttt{type},
\texttt{version}, \texttt{id}, \texttt{timestamp}, \texttt{cwd} --- and the timestamp comes
from the caller's clock, so a replay sees the same value. An existing file is left untouched:
a resumed session keeps its identity.

For \texttt{claude}, a fresh run names its session on the command line (\texttt{--session-id
<uuid>}), so a crash after the spawn still journals the handle. A resumed run carries
\texttt{--resume <sid>} instead.

The session id is the only handle that survives a turn cap or a killed process, and it is
unrecoverable once the process exits. That is why the dispatch scripts write it to a file at
start, and why the queue's recovery path reads the worker's own log when the dispatcher never
got to write it (chapter~\ref{ch:workflow}).

\paragraph{What the run reports about itself.} \texttt{reported\_routes} reads the agent's
\texttt{--mode json} output and collects the distinct \texttt{(provider, model)} pairs it
named. The distinction matters: a node asked for one route and the agent may have used
another, and an empty result means the output named no route at all --- which is reported as
unknown, never as a match. The same output yields the cost, the turn count and claude's own
error flag, all optional, because a killed run has an \texttt{init} line and no
\texttt{result} line.

\section{The role table that is specified and not landed}

Spec 362 exists because seven graphs hardcoded a model and three of them named an id that no
longer existed. The measured cost was concrete: the first real run of
\texttt{graphs/queue.toml} died in 1.6 seconds with

\begin{lstlisting}[language=bash]
{"code":"internal","exit_code":1,"message":"Model deepseek-v4-flash not found"}
\end{lstlisting}

and eight rows would have parked for one dead id. Two graphs already took the model from the
run input, which is why a stale id can hide for months: nothing validates it until an agent is
spawned.

The design is one config table per role and a role reference in a node:

\begin{lstlisting}[language=bash]
[models.workhorse]        # implements a spec
model = "deepseek-v4.1-flash"
provider = "ollama-cloud"

[models.planner]          # writes a spec
model = "glm-5.3"
provider = "ollama-cloud"
\end{lstlisting}

with R2 allowing \texttt{model = "@workhorse"} in a node, resolved when the node is built; R3
refusing an unknown role \emph{by name} and listing the configured roles, before any agent is
spawned; and R4 refusing an entry whose model is not on the provider's list when the provider
is reachable, while saying so and continuing when it is not --- an offline start must not
depend on the network.

R4 is the half that makes it real, and the spec says so in its own words: a worker can add the
table, resolve the role in one node kind, and leave every other call site hardcoded, so the
config exists, nothing uses it, and the next stale id hides in a graph nobody reads.

\begin{longtable}{@{}p{0.16\textwidth}p{0.26\textwidth}p{0.20\textwidth}p{0.30\textwidth}@{}}
\toprule
\textbf{Requirement} & \textbf{What it says} & \textbf{State at HEAD} & \textbf{Criterion} \\
\midrule
\endhead
R1 & one table per role, no model id in a graph & not landed & \texttt{grep workhorse} in
\texttt{config.rs} fails \\
R2 & node may name \texttt{@workhorse} & not landed & \texttt{grep '@workhorse' graphs/queue.toml}
fails \\
R3 & unknown role refused by name & not landed & \texttt{models\_from\_config} test file absent \\
R4 & dead model refused at config time & not landed & same test file \\
--- & no graph names a dead id & partially & \texttt{deepseek-v4-flash} still in four graph files \\
\bottomrule
\end{longtable}

The parked-backlog review of 2026-11-04 recorded the verdict as \texttt{needs-decision}, with
one open question for the operator: whether to raise the 450-effective-line budget for
\texttt{config.rs} or split the file further. The lesson recorded with it is the general one
--- the file-size guard counts effective lines, so a clean extraction that moves code without
deleting it stays over budget.

Until that is answered, the working rule is the one the tree enforces: a graph names its model
explicitly, and the daemon's \texttt{[core].model} is the fallback for graphs that take it
from input. A reader who wants a role table should treat the section above as the design and
check \texttt{graphs/queue.toml} for whether it has landed.

\section{Failure modes worth recognising}

\begin{enumerate}
	\item \textbf{A coder that reports success and wrote nothing.} Cause: no approval mode,
	every mutating tool denied, exit 3 with a plausible sentence. Check the argv before the
	report.
	\item \textbf{A coder that implemented another spec.} Cause: the spec path never reached
	the prompt. The prompt prefix is the fix; the symptom is a tree that matches the row's
	title and not its criteria.
	\item \textbf{A dispatch that dies in about a second with ``model not found''.} Cause: a
	stale id in a graph. This is the failure the role table is meant to move from spawn time to
	config time.
	\item \textbf{A dispatch refused with ``no credential for provider X''.} Cause: the
	provider is not in the two-entry key table, or its key is unset or empty in the daemon's
	environment. The message names the provider, so the fix is visible without reading code.
	\item \textbf{A run whose route is unknown.} Cause: text output, or output not captured.
	Unknown is reported as unknown; a missing route never reads as a match.
	\item \textbf{An escalation that never climbs.} Cause: the previous attempt failed without
	a failing criterion, so there is no evidence to escalate on. The ladder is doing what it
	says.
\end{enumerate}
